\clearpage
\oldpage[IV-3]{249}
\phantomsection
\addcontentsline{toc}{section}{Bibliografia (continuação)}
\section*{BIBLIOGRAFIA (\emph{continuação})}

\begingroup
\small
\setlength{\parindent}{0pt}
\setlength{\parskip}{3pt}
[39] J.-P.~\textsc{Serre}, \emph{Corps locaux}, Paris (Hermann), 1962.

[40] J.-P.~\textsc{Serre}, \emph{Groupes proalgébriques}, \emph{Publ. Inst. Hautes Études Sci.}, n\textsuperscript{o} 7 (1960).

\endgroup

\clearpage
\oldpage[IV-3]{250}
\phantomsection
\addcontentsline{toc}{section}{Índice de notações}
\section*{ÍNDICE DE NOTAÇÕES}

\begingroup
\small
\setlength{\parindent}{0pt}
\setlength{\parskip}{2pt}

$\mathfrak C(X)$, $\mathfrak D_c(X)$, $\mathfrak F_c(X)$,
$\mathfrak{LF}_c(X)$ ($X$ pré-esquema): 8.3.9.

$\mathfrak Q(\sh{F})$ ($\sh{F}$ $\sh{O}_X$-Módulo): 8.5.10.

$\mathfrak{Spr}(Y)$, $\mathfrak{Spro}(Y)$, $\mathfrak{Sprf}(Y)$
($Y$ pré-esquema): 8.6.1.

$\mathbf{Pro}(\mathcal C)$ ($\mathcal C$ categoria): 8.13.3.

$X_s$, $\sh{F}_s$, $u_s$, $g_s$, $Z_s$, $h_s$
($X$ $S$-pré-esquema, $s\in S$, $\sh{F}$ $\sh{O}_X$-Módulo,
$u$ homomorfismo de $\sh{O}_X$-Módulos, $g$ seção de
$\sh{O}_X$-Módulo, $Z$ parte de $X$, $h:X\to Y$ $S$-morfismo): 9.4.1.

$\mathfrak{Qc}(X)$, $\mathfrak{Lqc}(X)$, $\mathfrak O(X)$,
$\mathfrak F(X)$, $\mathfrak{Lf}(X)$, $\mathfrak{Lc}(X)$
($X$ espaço topológico): 10.1.1.

$\operatorname{prof}^{*}_{x}(\sh{F})$,
$\operatorname{prof}^{*}_{Z}(\sh{F})$,
$\operatorname{prof}^{*}(\sh{F})$
($\sh{F}$ $\sh{O}_X$-Módulo): 10.8.1.

$S(X)$ ($X$ pré-esquema de Jacobson): 10.9.1.

$S(f)$ ($f$ morfismo de pré-esquemas de Jacobson): 10.9.2.

$\operatorname{Spm}(A)$, $D^m(h)$ ($A$ anel de Jacobson, $h\in A$): 10.9.3.

$\operatorname{Tor.dim}_B(M)$ ($M$ $B$-módulo): 12.3.2.

\endgroup

\clearpage
\oldpage[IV-3]{251}
\phantomsection
\addcontentsline{toc}{section}{Índice terminológico}
\section*{ÍNDICE TERMINOLÓGICO}

\begingroup
\footnotesize
\setlength{\parindent}{0pt}
\setlength{\parskip}{1.1pt}
\newcommand{\EGAIVthreeTerm}[2]{#1: #2.\par}

\EGAIVthreeTerm{Aberto ultra-afim}{10.9.5}
\EGAIVthreeTerm{Critério de planicidade por fibras}{11.3.10}
\EGAIVthreeTerm{Espaço algébrico sobre $k$, espaço $k$-algébrico, espaço
pré-algébrico sobre $k$, espaço $k$-pré-algébrico}{10.10.2}
\EGAIVthreeTerm{Espaço de Jacobson}{10.3.1}
\EGAIVthreeTerm{Espaço $k$-pré-algébrico reduzido}{10.10.4}
\EGAIVthreeTerm{Espectro maximal de um anel de Jacobson}{10.9.3}
\EGAIVthreeTerm{Esqueleto primário de um Módulo, esqueleto virtual}{9.8.9}
\EGAIVthreeTerm{Família de morfismos esquematicamente dominante}{11.10.2}
\EGAIVthreeTerm{Família de morfismos universalmente esquematicamente dominante}{11.10.7}
\EGAIVthreeTerm{Família de subpré-esquemas esquematicamente densa}{11.10.2}
\EGAIVthreeTerm{Família separante de homomorfismos de feixes de módulos}{11.9.1}
\EGAIVthreeTerm{Família separante de homomorfismos de módulos}{11.9.4}
\EGAIVthreeTerm{Família universalmente separante de homomorfismos de feixes de módulos}{11.9.14}
\EGAIVthreeTerm{Geometricamente isomorfos ($k$-pré-esquemas)}{9.1.4}
\EGAIVthreeTerm{Homomorfismo de Módulos universalmente injetivo}{11.9.14 e 11.9.18}
\EGAIVthreeTerm{Lema de Chow para os morfismos de apresentação finita}{8.10.5.1}
\EGAIVthreeTerm{Morfismo equidimensional em um ponto, morfismo equidimensional}{13.2.2 e 13.3.2}
\EGAIVthreeTerm{Morfismo próprio em um ponto}{15.7.1}
\EGAIVthreeTerm{Morfismo pseudofinito}{8.12.3}
\EGAIVthreeTerm{Morfismo submersivo, morfismo universalmente submersivo}{15.7.8}
\EGAIVthreeTerm{Morfismo universalmente aberto em um ponto}{14.3.3}
\EGAIVthreeTerm{Parte finita saturada de um $k$-pré-esquema algébrico}{9.8.9}
\EGAIVthreeTerm{Parte muito densa de um espaço topológico}{10.1.3}
\EGAIVthreeTerm{Parte própria em um ponto (de um $S$-pré-esquema)}{15.7.1}
\EGAIVthreeTerm{Parte quase-construtível, parte localmente quase-construtível de um espaço topológico}{10.1.1}
\EGAIVthreeTerm{Polinômio geometricamente irredutível}{9.7.4}
\EGAIVthreeTerm{Pré-esquema de Jacobson}{10.4.1}
\EGAIVthreeTerm{Pré-esquema equidimensional sobre outro em um ponto, equidimensional sobre outro}{13.2.2 e 13.3.2}
\EGAIVthreeTerm{Pré-esquema essencialmente afim sobre outro}{8.13.4}
\EGAIVthreeTerm{Princípio da extensão finita}{9.1.1}
\EGAIVthreeTerm{Profundidade retificada}{10.8.1}
\EGAIVthreeTerm{Pro-objeto, morfismo de pro-objetos}{8.13.3}
\EGAIVthreeTerm{Pro-objeto essencialmente afim}{8.13.4, 8.14.1}
\EGAIVthreeTerm{Propriedade construtível, propriedade ind-construtível}{9.2.1 e 9.2.2, (i) e (ii)}
\EGAIVthreeTerm{Quase-homeomorfismo de espaços topológicos}{10.2.2}
\EGAIVthreeTerm{Quase-isomorfismo de espaços anelados}{10.2.8}
\EGAIVthreeTerm{$R$-ultrapré-esquema}{10.10.2}
\EGAIVthreeTerm{Saturada de uma parte finita de um $k$-pré-esquema algébrico}{9.8.9}
\EGAIVthreeTerm{Subconjunto de um $S$-pré-esquema submersivo, universalmente submersivo sobre sua imagem}{15.7.8}
\EGAIVthreeTerm{``Teorema principal'' de Zariski para os morfismos de apresentação finita}{8.12.6}
\EGAIVthreeTerm{Tipo primário de um Módulo, tipo primário geométrico de um Módulo}{9.8.9}
\EGAIVthreeTerm{Ultrapré-esquema, morfismo de ultrapré-esquemas}{10.9.5}

\let\EGAIVthreeTerm\relax
\endgroup

\clearpage
% A página impressa 252 está em branco.
\thispagestyle{empty}
\null

\clearpage
\oldpage[IV-3]{253}
\phantomsection
\addcontentsline{toc}{section}{Sumário}
\section*{SUMÁRIO}

\begingroup
\small
\setlength{\parindent}{0pt}
\setlength{\parskip}{1.5pt}
\newcommand{\EGAIVthreeContentsLine}[3]{\textbf{#1}\quad #2\unskip\nobreak\hspace{0.4em}\nobreak\dotfill\nobreak\mbox{#3}\par}

\EGAIVthreeContentsLine{CAPÍTULO IV.}{\textbf{Estudo local dos esquemas e dos morfismos de esquemas} \emph{(continuação)}}{5}
\EGAIVthreeContentsLine{\S~8.}{Limites projetivos de pré-esquemas}{5}
\EGAIVthreeContentsLine{8.1.}{Introdução}{5}
\EGAIVthreeContentsLine{8.2.}{Limites projetivos de pré-esquemas}{7}
\EGAIVthreeContentsLine{8.3.}{Subconjuntos construtíveis de um limite projetivo de pré-esquemas}{12}
\EGAIVthreeContentsLine{8.4.}{Critérios de irredutibilidade e conexidade para os limites projetivos de pré-esquemas}{17}
\EGAIVthreeContentsLine{8.5.}{Módulos de apresentação finita sobre um limite projetivo de pré-esquemas}{19}
\EGAIVthreeContentsLine{8.6.}{Subpré-esquemas de apresentação finita de um limite projetivo de pré-esquemas}{25}
\EGAIVthreeContentsLine{8.7.}{Critérios para que um limite projetivo de pré-esquemas seja um pré-esquema reduzido (resp.\ integral)}{27}
\EGAIVthreeContentsLine{8.8.}{Pré-esquemas de apresentação finita sobre um limite projetivo de pré-esquemas}{28}
\EGAIVthreeContentsLine{8.9.}{Primeiras aplicações à eliminação das hipóteses noetherianas}{34}
\EGAIVthreeContentsLine{8.10.}{Propriedades de permanência dos morfismos por passagem ao limite projetivo}{36}
\EGAIVthreeContentsLine{8.11.}{Aplicação aos morfismos quase-finitos}{41}
\EGAIVthreeContentsLine{8.12.}{Nova demonstração e generalização do ``Teorema principal'' de Zariski}{43}
\EGAIVthreeContentsLine{8.13.}{Tradução em termos de pro-objetos}{49}
\EGAIVthreeContentsLine{8.14.}{Caracterização de um pré-esquema localmente de apresentação finita sobre outro em termos do funtor que ele representa}{52}

\EGAIVthreeContentsLine{\S~9.}{Propriedades construtíveis}{54}
\EGAIVthreeContentsLine{9.1.}{O princípio da extensão finita}{54}
\EGAIVthreeContentsLine{9.2.}{Propriedades construtíveis e ind-construtíveis}{56}
\EGAIVthreeContentsLine{9.3.}{Propriedades construtíveis dos morfismos de pré-esquemas algébricos}{60}

\clearpage
\oldpage[IV-3]{254}
\EGAIVthreeContentsLine{9.4.}{Construtibilidade de certas propriedades dos módulos}{62}
\EGAIVthreeContentsLine{9.5.}{Construtibilidade de propriedades topológicas}{67}
\EGAIVthreeContentsLine{9.6.}{Construtibilidade de certas propriedades de morfismos}{71}
\EGAIVthreeContentsLine{9.7.}{Construtibilidade das propriedades de separabilidade, de irredutibilidade geométrica e de conexidade geométrica}{76}
\EGAIVthreeContentsLine{9.8.}{Decomposição primária em uma vizinhança de uma fibra genérica}{83}
\EGAIVthreeContentsLine{9.9.}{Construtibilidade das propriedades locais das fibras}{88}

\EGAIVthreeContentsLine{\S~10.}{Pré-esquemas de Jacobson}{95}
\EGAIVthreeContentsLine{10.1.}{Subconjuntos muito densos de um espaço topológico}{95}
\EGAIVthreeContentsLine{10.2.}{Quase-homeomorfismos}{97}
\EGAIVthreeContentsLine{10.3.}{Espaços de Jacobson}{101}
\EGAIVthreeContentsLine{10.4.}{Pré-esquemas de Jacobson e anéis de Jacobson}{101}
\EGAIVthreeContentsLine{10.5.}{Pré-esquemas de Jacobson noetherianos}{104}
\EGAIVthreeContentsLine{10.6.}{Dimensão nos pré-esquemas de Jacobson}{107}
\EGAIVthreeContentsLine{10.7.}{Exemplos e contraexemplos}{109}
\EGAIVthreeContentsLine{10.8.}{Profundidade retificada}{110}
\EGAIVthreeContentsLine{10.9.}{Espectros maximais e ultrapré-esquemas}{112}
\EGAIVthreeContentsLine{10.10.}{Espaços algébricos de Serre}{114}

\EGAIVthreeContentsLine{\S~11.}{Propriedades topológicas dos morfismos planos de apresentação finita. Critérios de planicidade}{116}
\EGAIVthreeContentsLine{11.1.}{Conjuntos de planicidade (caso noetheriano)}{117}
\EGAIVthreeContentsLine{11.2.}{Planicidade de um limite projetivo de pré-esquemas}{119}
\EGAIVthreeContentsLine{11.3.}{Aplicação à eliminação de hipóteses noetherianas}{132}
\EGAIVthreeContentsLine{11.4.}{Descida da planicidade por morfismos quaisquer: caso de um pré-esquema de base artiniano}{143}
\EGAIVthreeContentsLine{11.5.}{Descida da planicidade por morfismos quaisquer: caso geral}{150}
\EGAIVthreeContentsLine{11.6.}{Descida da planicidade por morfismos quaisquer: caso de um pré-esquema de base unirramificado}{154}
\EGAIVthreeContentsLine{11.7.}{Contraexemplos}{157}
\EGAIVthreeContentsLine{11.8.}{Um critério valuativo de planicidade}{159}
\EGAIVthreeContentsLine{11.9.}{Famílias separantes e universalmente separantes de homomorfismos de feixes de módulos}{160}
\EGAIVthreeContentsLine{11.10.}{Famílias esquematicamente dominantes de morfismos e famílias esquematicamente densas de subpré-esquemas}{170}

\EGAIVthreeContentsLine{\S~12.}{Estudo das fibras dos morfismos planos de apresentação finita}{173}
\EGAIVthreeContentsLine{12.0.}{Introdução}{173}
\EGAIVthreeContentsLine{12.1.}{Propriedades locais das fibras de um morfismo plano localmente de apresentação finita}{174}

\clearpage
\oldpage[IV-3]{255}
\EGAIVthreeContentsLine{12.2.}{Propriedades locais e globais das fibras de um morfismo próprio, plano e de apresentação finita}{179}
\EGAIVthreeContentsLine{12.3.}{Propriedades cohomológicas locais das fibras de um morfismo plano e localmente de apresentação finita}{183}

\EGAIVthreeContentsLine{\S~13.}{Morfismos equidimensionais}{187}
\EGAIVthreeContentsLine{13.1.}{O teorema de semicontinuidade de Chevalley}{188}
\EGAIVthreeContentsLine{13.2.}{Morfismos equidimensionais: caso dos morfismos dominantes de pré-esquemas irredutíveis}{190}
\EGAIVthreeContentsLine{13.3.}{Morfismos equidimensionais: caso geral}{194}

\EGAIVthreeContentsLine{\S~14.}{Morfismos universalmente abertos}{199}
\EGAIVthreeContentsLine{14.1.}{Morfismos abertos}{200}
\EGAIVthreeContentsLine{14.2.}{Morfismos abertos e fórmula da dimensão}{202}
\EGAIVthreeContentsLine{14.3.}{Morfismos universalmente abertos}{204}
\EGAIVthreeContentsLine{14.4.}{Critério de Chevalley para os morfismos universalmente abertos}{209}
\EGAIVthreeContentsLine{14.5.}{Morfismos universalmente abertos e quase-seções}{216}

\EGAIVthreeContentsLine{\S~15.}{Estudo das fibras de um morfismo universalmente aberto}{223}
\EGAIVthreeContentsLine{15.1.}{Multiplicidades das fibras de um morfismo universalmente aberto}{223}
\EGAIVthreeContentsLine{15.2.}{Planicidade dos morfismos universalmente abertos com fibras geometricamente reduzidas}{226}
\EGAIVthreeContentsLine{15.3.}{Aplicações: critérios de redução e irredutibilidade}{228}
\EGAIVthreeContentsLine{15.4.}{Complementos sobre os morfismos de Cohen--Macaulay}{229}
\EGAIVthreeContentsLine{15.5.}{Posto separável das fibras de um morfismo quase-finito e universalmente aberto. Aplicação às componentes conexas geométricas das fibras de um morfismo próprio}{231}
\EGAIVthreeContentsLine{15.6.}{Componentes conexas das fibras ao longo de uma seção}{236}
\EGAIVthreeContentsLine{15.7.}{Apêndice: critérios valuativos de propriedade local}{242}
\EGAIVthreeContentsLine{}{BIBLIOGRAFIA \emph{(continuação)}}{249}
\EGAIVthreeContentsLine{}{ÍNDICE DE NOTAÇÕES}{250}
\EGAIVthreeContentsLine{}{ÍNDICE TERMINOLÓGICO}{251}

\vspace{2em}
\begin{center}
\emph{Manuscrito recebido em 15 de novembro de 1964.}
\end{center}

\let\EGAIVthreeContentsLine\relax
\endgroup
